% TO DO
% Italizar
% Revisar nuevamente las fuentes 
% eso de paradigma de VA no m termina de convencer
% Aporte computacional en la introducción — reforzar que las preguntas de investigación solo pueden responderse construyendo y evaluando el sistema, para que los objetivos se lean como investigación y no como ingeniería.
%Poner bonito la organizacion d la tesis
\chapter{Introducción}
\label{cap:introduccion}

Cada año, instituciones públicas como tribunales, organismos de 
control y agencias regulatorias producen miles de documentos que 
registran el avance de procesos que pueden extenderse durante meses 
o años. Identificar qué ocurrió en un caso, cuándo y dónde, sigue 
requiriendo que un analista lea cientos de documentos de forma 
manual, porque la información está dispersa en archivos generados 
por distintos actores en distintos momentos. Reconstruir 
automáticamente la secuencia de eventos que conforma cada caso 
sigue siendo un problema abierto. Los sistemas tradicionales de 
extracción de información operan a nivel de documento individual y 
los enfoques basados en reglas requieren plantillas específicas por 
tipo de documento, lo que dificulta integrar información cuando una 
misma entidad aparece referenciada de formas distintas en archivos 
generados por distintos actores~\cite{cattan2021crossdoc}.

Los Grandes Modelos de Lenguaje (\acs{LLM}, por sus siglas en inglés)
 han demostrado capacidades 
prometedoras para extraer entidades, relaciones y secuencias 
narrativas de documentos técnicos con una flexibilidad que los 
enfoques basados en reglas no alcanzan~\cite{brown2020, wei2022}. 
Sin embargo, en tareas de extracción estructurada estos modelos 
presentan tasas de alucinación significativas~\cite{ji2023}, lo 
cual representa un riesgo inaceptable cuando una fecha incorrecta 
o un monto fabricado puede comprometer una decisión institucional. 
Si bien parte de estas inconsistencias puede detectarse 
automáticamente contrastando cada dato contra el texto fuente, la 
corrección automática no resuelve el problema de fondo. El analista 
necesita un marco que le permita calibrar su confianza en lo 
extraído mientras explora cómo evolucionó el proceso. El 
\acl{VA} (\acs{VA}) ~\cite{keim2008visual, sacha2014} 
ofrece ese marco. A diferencia de las interfaces de consulta 
basadas en chat, que devuelven respuestas puntuales sobre un 
documento a la vez, %Creo que aqui falta una cita :)
\ac{VA} permite mantener simultáneamente las 
dimensiones temporal, geoespacial y semántica de un proceso 
distribuido en múltiples documentos, integrando el procesamiento 
automático con el razonamiento humano. Esta integración aún no ha 
sido evaluada en documentos institucionales fragmentados en español.

En este trabajo abordamos esa brecha mediante el diseño y 
evaluación de un sistema de \ac{VA} que articula tres componentes: un 
pipeline de extracción basado en \acs{LLM} con corrección automática de 
inconsistencias, un mecanismo de trazabilidad que vincula cada 
dato con su fragmento de origen, y un conjunto de vistas 
interactivas coordinadas para las dimensiones temporal, geoespacial 
y narrativa de los procesos. Concretamente, esta investigación 
responde dos preguntas. \textbf{(1)} ¿Cómo pueden los \acs{LLM} automatizar la extracción de información estructurada a partir de documentos institucionales fragmentados en español? \textbf{(2)} 
¿Qué formas de visualización e interacción permiten a un 
analista explorar y verificar narrativas complejas extraídas 
de documentos institucionales fragmentados? Para 
responderlas, evaluamos el sistema con expedientes reales de 
auditoría de la Contraloría General de la República del Perú, un contexto real con alto volumen documental en español.




\section{Motivación y Contexto}

Los organismos de control gubernamental generan grandes 
volúmenes de documentos que registran procesos extendidos en 
el tiempo. En el Perú, la Contraloría General de la República 
procesó más de 50,000 informes solo en 
2024~\cite{andina2024}. Estos documentos presentan 
características que los convierten en un caso de estudio 
relevante para la investigación en análisis visual de texto: 
combinan narrativa densa con terminología especializada, 
involucran múltiples actores y ubicaciones geográficas, y 
registran secuencias de eventos que se extienden durante 
meses a través de documentos generados de forma 
independiente.

Si bien el \ac{NLP} permite clasificar y extraer elementos 
clave de documentos legales~\cite{zhong-etal-2020-nlp, 
lippi2019legal}, estos enfoques operan predominantemente a 
nivel de documento individual y requieren adaptación 
específica por dominio. Los \acs{LLM}, en cambio, permiten 
adaptar la extracción a nuevos tipos de documentos sin 
reentrenamiento, lo que los convierte en una alternativa 
de propósito general para múltiples dominios.

El \ac{VA} aborda este tipo de problemas combinando procesamiento automático con representación visual interactiva, permitiendo que el analista explore grandes volúmenes de información mientras mantiene control sobre el proceso de interpretación~\cite{keim2008visual, sacha2014}. Esta combinación es especialmente relevante cuando los datos extraídos automáticamente requieren validación experta antes de sustentar decisiones.

Sin embargo, reconstruir la secuencia completa de un proceso 
a partir de múltiples documentos fragmentados requiere 
resolver desafíos adicionales tales como identificar que 
distintas menciones en distintos archivos se refieren a la 
misma entidad~\cite{cattan2021crossdoc}, ordenar 
temporalmente eventos descritos con formatos heterogéneos, 
y consolidar información dispersa en una representación 
coherente. Permitir además que el analista verifique esa 
reconstrucción sigue siendo un problema 
abierto~\cite{pereira2023inacia}. Abordar este problema de 
forma integrada, combinando extracción automática con 
representación visual interactiva, es el punto de partida 
de esta investigación.

\section{Planteamiento del Problema}

El análisis de procesos documentados en texto no estructurado 
plantea tres problemas computacionales sin resolver.

El primero es la extracción y resolución de entidades a través 
de múltiples documentos. Los sistemas existentes operan 
predominantemente a nivel de documento individual y no resuelven 
la correferencia entre documentos: la misma entidad puede aparecer 
referenciada de formas distintas en distintos archivos, y sin un 
mecanismo de resolución de identidad los datos extraídos no pueden 
integrarse de manera coherente. Los \acs{LLM} ofrecen capacidades 
prometedoras para este problema, pero su comportamiento en 
documentos técnicos y jurídicos en español no ha sido 
sistemáticamente evaluado en pipelines de extracción con revisión 
asistida~\cite{zhong-etal-2020-nlp}.

El segundo es la verificabilidad de la extracción automática. 
Cuando un sistema extrae un dato, el analista no dispone de 
mecanismos inmediatos para verificar su corrección sin consultar 
manualmente el documento fuente. En dominios donde un dato incorrecto puede sustentar una decisión institucional, esta ausencia de verificación hace que los sistemas automáticos sean difíciles de adoptar en la práctica~\cite{ji2023}.

El tercero es la representación visual integrada de procesos 
multidimensionales. Los procesos documentados en texto tienen 
simultáneamente una dimensión temporal, espacial y narrativa. 
Caracterizar qué representaciones visuales e interacciones permiten explorar de forma integrada las dimensiones temporal, espacial y narrativa de un proceso distribuido en múltiples documentos sigue siendo un problema abierto~\cite{zhao2023contextwing, rajabiyazdi2024textvista, ye2024storyexplorer, yeh2025storyribbons}.
\section{Objetivos}

\subsection{Objetivo General}

Desarrollar un sistema de \acl{VA} con extracción automática de 
información mediante \acs{LLM} y revisión asistida, para la 
exploración interactiva de procesos documentados en texto 
no estructurado a través de sus dimensiones temporal, 
geoespacial y narrativa.

\subsection{Objetivos Específicos}

\begin{enumerate}
    \item Diseñar un \textit{pipeline} de extracción basado en \acs{LLM} 
     que identifique entidades, hitos temporales, actores, 
     relaciones y ubicaciones en documentos textuales, 
     mitigando posibles alucinaciones del modelo.

     % ambas vistas son formas complementarias de explorar los datos extraídos, por eso están en el mismo objetivo
     \item Desarrollar un prototipo de vistas interactivas 
     coordinadas que permitan explorar las dimensiones 
     temporal, geoespacial y narrativa de los procesos 
     documentados, incluyendo una vista de similitud 
     semántica que permita comparar entidades entre sí.
     
     %\item Desarrollar un mecanismo de trazabilidad que preserve la referencia al documento fuente durante todo el pipeline, permitiendo al analista acceder al fragmento original de cada dato desde la interfaz.
     
     \item Evaluar la precisión del pipeline de extracción 
     mediante un caso de estudio con expedientes de la 
     Contraloría General de la República del Perú, 
     comparando las extracciones generadas contra el 
     contenido verificado de los documentos fuente.
     
     \item Evaluar la utilidad y usabilidad del sistema 
     mediante una prueba con usuarios, midiendo eficiencia 
     en tareas de exploración documental y utilidad 
     percibida mediante cuestionarios estandarizados.
\end{enumerate}


\section{Organización de la tesis}

El presente documento se organiza de la siguiente manera:

\begin{itemize}
    \item \hyperref[cap:introduccion]{\textbf{Capítulo~\ref*{cap:introduccion}: Introducción}} presenta el problema de investigación, la motivación, las preguntas de investigación y los objetivos del trabajo.
    \item \hyperref[cap:marco_teorico]{\textbf{Capítulo~\ref*{cap:marco_teorico}: Marco Teórico}} desarrolla el marco teórico, 
     cubriendo los fundamentos de \ac{VA}, visualización 
     espacio-temporal, extracción de información con \acs{LLM} 
      y el dominio de Control Concurrente de la Contraloría 
      General de la República del Perú.
     \item \hyperref[cap:estado_arte]{\textbf{Capítulo~\ref*{cap:estado_arte}: Estado del Arte}} presenta el estado del arte, 
      revisando los trabajos relacionados en extracción de 
      información con \acs{LLM}, sistemas de \ac{VA} para texto y 
      visualización de narrativas complejas.
      \item \hyperref[cap:propuesta]{\textbf{Capítulo~\ref*{cap:propuesta}: Propuesta}} describe el diseño e implementación 
      del sistema, incluyendo el pipeline de extracción, la 
      corrección automática de inconsistencias y las vistas 
      interactivas coordinadas.
      \item El Capítulo 5 presenta la evaluación del sistema 
      en tres partes: la precisión del pipeline y el impacto 
      de la corrección automática sobre la tasa de 
      alucinaciones, la calidad de la proyección semántica 
      para identificación de patrones, y la prueba de 
      usabilidad con usuarios.
      \item El Capítulo 6 discute las conclusiones, las 
      limitaciones del trabajo y las líneas de investigación 
      futura.
\end{itemize}

\section{Cronograma}

\begin{figure}[H]
    \centering
    \includegraphics[width=0.9\textwidth]{Tesis_CS_202203/Cro_Actv.pdf}
    \caption{Cronograma de actividades.}
    \label{fig:cronograma}
\end{figure}
